% Flush all pending floats from the experiments section so that no table or figure
% is deferred past the start of the conclusion (avoids the stranded-conclusion and
% end-of-paper float-interruption issue).
\clearpage
\section{Limitations and Conclusion}

\paragraph{Limitations.}
\method{} inherits three scope boundaries. It depends on local replayability
(AS2): fully non-replayable, single-shot environments cannot yield \rit{} labels,
and world-model-simulated replay or off-policy evaluation would be needed as a
substitute. It is subject to an interaction caveat (AS3): leave-one-out does not
capture higher-order interactions, most visibly for redundant/substitutable
memories, and the exploratory group-intervention analysis that surfaces them adds
cost and carries no numerical guarantee. And
its benefit has a boundary (H5): single-domain short streams and small banks are
expected to yield little gain, so \method{} is intended for long-term deployment
rather than short sessions. We further scope the contribution to credit assignment
for existing memories and leave write-time quality out of scope; \method{}
composes with methods that address it. We emphasize that governance and scope
gating are \emph{empirically tested consumers} of the immediate-credit signal, not
validated causal estimators of long-horizon retention value. A rebuttal-sensitive
risk is that, if
attribution scores can be influenced by externally supplied content, the audit
becomes a new attack surface for memory-poisoning ``credit laundering,'' which we
flag for future defense. Whether \aca{} transfers across fundamentally different
modalities beyond the tested text and agentic environments is not specified by our
source and is left as an open question, excluded from C1--C4.

\paragraph{Conclusion.}
We argued that self-evolving agent memory is limited, in its credit-assignment
step, less by better heuristics than by a missing primitive: reliable immediate
counterfactual credit, in the immediate-contribution setting we study.
Correlational bookkeeping mistakes retrieval for contribution and, we hypothesize,
is associated with superstitious memories that accumulate over deployment.
\method{} is designed to supply this counterfactual credit
affordably by collecting small-budget interventional labels (\rit{}), amortizing
them into a per-event score that is constant in bank size (\aca{}), and exposing
that single score to two
minimal consumers. The evaluation is planned rather than completed: it tests
whether \aca{} recovers the audited signal in closed loop, whether its consumers
flatten superstition and reduce cross-task interference, and whether the recovered
signal improves average success at low overhead---each behind an explicit
falsification gate. Two concrete directions follow: structure-aware interaction
attribution to relax AS3, and replay substitutes to relax AS2.
